\documentclass{article}
\usepackage{graphicx}
\usepackage{amsmath,amssymb,amsthm,amsfonts}
\usepackage{booktabs}
\usepackage[authoryear,round]{natbib}
\usepackage[hidelinks]{hyperref}

\title{Supplementary Material: Mixtures of Generalized Regression Models with
Component-Specific Response Families}
\author{Samyajoy Pal\textsuperscript{1}\quad
Dayasri Ravi\textsuperscript{2}\quad
Rouven E. Haschka\textsuperscript{1}\\[0.75em]
\small \textsuperscript{1}Chair of Data Analytics, Technical University of Kaiserslautern-Landau\\
\small \textsuperscript{2}Department of Mathematics, University of Augsburg}
\date{}

\newtheorem{remark}{Remark}

\DeclareMathOperator{\logit}{logit}
\newcommand{\E}{\mathbb{E}}
\newcommand{\Var}{\operatorname{Var}}
\newcommand{\R}{\mathbb{R}}
\setlength{\emergencystretch}{2em}
\renewcommand{\thesection}{S\arabic{section}}
\renewcommand{\thesubsection}{S\arabic{section}.\arabic{subsection}}
\renewcommand{\thetable}{S\arabic{table}}
\renewcommand{\thefigure}{S\arabic{figure}}

\begin{document}

\maketitle

\section{Family-Specific Analytic Louis Derivatives}
\label{sec:supp_louis}

The main manuscript states the finite-mixture Louis identity, the
regression-coordinate chain rule, and the family-specific derivative blocks in
Appendix A.  This supplement is complementary: it records the algebra in a
more expansive form, documents implementation validation, and supplies the
simulation and application details that would interrupt the flow of the main
paper.

The derivations below cover the families with analytic derivative blocks:
Gaussian, Student-$t$, Poisson, negative binomial (NB2), Gamma, exponential,
lognormal, inverse Gaussian, Bernoulli, geometric, beta, zero-inflated Poisson
(ZIP), zero-inflated negative binomial (ZINB), and skew normal.  Families whose
densities are delegated to external special-function implementations are not
claimed to have closed-form Louis blocks in the present implementation.

\subsection{Notation}

For one component and one observation, write
\[
    a_\mu=\frac{\partial \ell}{\partial\mu},\qquad
    b_{\mu\mu}=\frac{\partial^2\ell}{\partial\mu^2},\qquad
    a_\gamma=\frac{\partial\ell}{\partial\gamma},\qquad
    B_{\gamma\gamma}=\frac{\partial^2\ell}{\partial\gamma\,\partial\gamma^\top},
    \qquad
    b_{\mu\gamma}=\frac{\partial^2\ell}{\partial\mu\,\partial\gamma}.
\]
Let $\mu=m(\eta)$, $\dot m(\eta)=\partial m(\eta)/\partial\eta$, and
$\ddot m(\eta)=\partial^2m(\eta)/\partial\eta^2$.  The main manuscript's
regression-coordinate proposition maps the blocks below into
$(\beta,\gamma)$ coordinates using the generalized linear model (GLM) link.

For the links used in the simulations,
\[
    \begin{gathered}
    \text{identity:}\quad m(\eta)=\eta,\quad
    \dot m=1,\quad \ddot m=0,\\
    \text{log:}\quad m(\eta)=\exp(\eta),\quad
    \dot m=m,\quad \ddot m=m,\\
    \text{logit:}\quad m(\eta)=\{1+\exp(-\eta)\}^{-1},\\
    \dot m=m(1-m),\quad
    \ddot m=m(1-m)(1-2m).
    \end{gathered}
\]

\subsection{Analytic blocks by family}

All formulas below are for a single observation.  Derivatives are in
$(\mu,\gamma)$ coordinates, where $\gamma$ denotes the transformed nuisance
coordinate used by the optimizer.  Cross-derivatives not displayed are zero.
The functions $\psi$ and $\psi_1$ denote the digamma and trigamma functions.

\subsubsection*{Gaussian}

For $Y\sim N(\mu,\sigma^2)$, let $v=\sigma^2=\exp(s)$ with
$s=\log\sigma^2$ and $r=y-\mu$.  Then
\[
    a_\mu=\frac{r}{v},\qquad
    b_{\mu\mu}=-\frac{1}{v},
\]
\[
    a_s=-\frac12+\frac{r^2}{2v},\qquad
    B_{ss}=-\frac{r^2}{2v},\qquad
    b_{\mu s}=-\frac{r}{v}.
\]

\subsubsection*{Student-$t$}

For a Student-$t$ component with location $\mu$, scale
$\sigma=\exp(a)$, and degrees of freedom $\nu=2+\exp(b)$, define
$w=\exp(b)$, $r=y-\mu$, $D=\nu\sigma^2+r^2$, $u=r^2/\sigma^2$, and
$q=1+u/\nu$.  Then
\[
    a_\mu=\frac{(\nu+1)r}{D},\qquad
    b_{\mu\mu}=\frac{(\nu+1)(r^2-\nu\sigma^2)}{D^2},
\]
\[
    a_a=-1+\frac{(\nu+1)r^2}{D},\qquad
    B_{aa}=-\frac{2(\nu+1)r^2\nu\sigma^2}{D^2},\qquad
    b_{\mu a}=-\frac{2(\nu+1)r\nu\sigma^2}{D^2}.
\]
Let
\[
    L_\nu=
    \frac12\psi\!\left(\frac{\nu+1}{2}\right)
    -\frac12\psi\!\left(\frac{\nu}{2}\right)
    -\frac{1}{2\nu}
    -\frac12\log q
    +\frac{(\nu+1)u}{2\nu(\nu+u)} .
\]
Also define
\[
    M=\nu(\nu+u),\qquad
    N=\frac12u(\nu+1),
\]
\[
    L_{\nu\nu}
    =
    \frac14\psi_1\!\left(\frac{\nu+1}{2}\right)
    -\frac14\psi_1\!\left(\frac{\nu}{2}\right)
    +\frac{1}{2\nu^2}
    +\frac{u}{2M}
    +\frac{\frac12uM-N(2\nu+u)}{M^2}.
\]
The transformed degree-of-freedom derivatives are
\[
    a_b=wL_\nu,\qquad
    B_{bb}=wL_\nu+w^2L_{\nu\nu},
\]
\[
    b_{\mu b}
    =
    \frac{w\,r(r^2-\sigma^2)}{D^2},
    \qquad
    B_{ab}
    =
    \frac{w\,r^2(r^2-\sigma^2)}{D^2}.
\]

\subsubsection*{Poisson}

For $Y\sim\operatorname{Poisson}(\mu)$,
\[
    a_\mu=\frac{y}{\mu}-1,\qquad
    b_{\mu\mu}=-\frac{y}{\mu^2}.
\]

\subsubsection*{Negative binomial NB2}

For the NB2 parameterization $\Var(Y)=\mu+\alpha\mu^2$, let
$\rho=\log\alpha$, $r=\alpha^{-1}=\exp(-\rho)$, and $A=r+\mu$.  The log pmf is
\[
    \ell=
    \log\Gamma(y+r)-\log\Gamma(r)-\log\Gamma(y+1)
    +r\log r+y\log\mu-(r+y)\log(r+\mu).
\]
Then
\[
    a_\mu=\frac{y}{\mu}-\frac{r+y}{A},
    \qquad
    b_{\mu\mu}=-\frac{y}{\mu^2}+\frac{r+y}{A^2}.
\]
Let
\[
    L_r=\psi(y+r)-\psi(r)+\log r+1-\log A-\frac{r+y}{A},
\]
\[
    L_{rr}
    =
    \psi_1(y+r)-\psi_1(r)+\frac1r-\frac1A-\frac{\mu-y}{A^2}.
\]
Since $\partial r/\partial\rho=-r$,
\[
    a_\rho=-rL_r,\qquad
    B_{\rho\rho}=rL_r+r^2L_{rr},\qquad
    b_{\mu\rho}=\frac{r(\mu-y)}{A^2}.
\]

\subsubsection*{Gamma}

For the mean-shape Gamma parameterization with
$\E(Y)=\mu$ and shape $\kappa=\exp(a)$,
\[
    \ell
    =
    (\kappa-1)\log y-\kappa y/\mu-\log\Gamma(\kappa)
    -\kappa\log\mu+\kappa\log\kappa .
\]
Define
\[
    L_\kappa=\log y-y/\mu-\psi(\kappa)-\log\mu+\log\kappa+1,
    \qquad
    L_{\kappa\kappa}=-\psi_1(\kappa)+\frac1\kappa.
\]
Then
\[
    a_\mu=\frac{\kappa(y-\mu)}{\mu^2},\qquad
    b_{\mu\mu}=\frac{\kappa(\mu-2y)}{\mu^3},
\]
\[
    a_a=\kappa L_\kappa,\qquad
    B_{aa}=\kappa L_\kappa+\kappa^2L_{\kappa\kappa},\qquad
    b_{\mu a}=a_\mu.
\]

\subsubsection*{Exponential}

For the exponential distribution parameterized by mean $\mu$,
\[
    \ell=-\log\mu-y/\mu,\qquad
    a_\mu=\frac{y-\mu}{\mu^2},\qquad
    b_{\mu\mu}=\frac{\mu-2y}{\mu^3}.
\]

\subsubsection*{Lognormal}

For the lognormal component, $\mu$ denotes the log-location, not the response
mean.  Let $\sigma=\exp(a)$ and $r=\log y-\mu$.  Then
\[
    a_\mu=\frac{r}{\sigma^2},\qquad
    b_{\mu\mu}=-\frac{1}{\sigma^2},
\]
\[
    a_a=-1+\frac{r^2}{\sigma^2},\qquad
    B_{aa}=-\frac{2r^2}{\sigma^2},\qquad
    b_{\mu a}=-\frac{2r}{\sigma^2}.
\]

\subsubsection*{Inverse Gaussian}

For the inverse Gaussian density with mean $\mu$ and
$\lambda=\exp(a)$, define
\[
    A=\frac{(y-\mu)^2}{2\mu^2y}.
\]
Then
\[
    a_\mu=\frac{\lambda(y-\mu)}{\mu^3},\qquad
    b_{\mu\mu}=\frac{\lambda(2\mu-3y)}{\mu^4},
\]
\[
    a_a=\frac12-\lambda A,\qquad
    B_{aa}=-\lambda A,\qquad
    b_{\mu a}=a_\mu.
\]

\subsubsection*{Bernoulli}

For $Y\in\{0,1\}$ with success probability $\mu$,
\[
    a_\mu=\frac{y}{\mu}-\frac{1-y}{1-\mu},\qquad
    b_{\mu\mu}=-\frac{y}{\mu^2}-\frac{1-y}{(1-\mu)^2}.
\]

\subsubsection*{Geometric}

For the geometric distribution on $\{0,1,\ldots\}$ parameterized by its mean
$\mu$,
\[
    \ell=y\log\mu-(y+1)\log(1+\mu),
\]
and therefore
\[
    a_\mu=\frac{y}{\mu}-\frac{y+1}{1+\mu},\qquad
    b_{\mu\mu}=-\frac{y}{\mu^2}+\frac{y+1}{(1+\mu)^2}.
\]

\subsubsection*{Beta}

For the beta mean-precision parameterization, let
$\phi=\exp(a)$, $A=\mu\phi$, and $B=(1-\mu)\phi$.  Define
\[
    G=-\psi(A)+\psi(B)+\log y-\log(1-y),
\]
\[
    L_\phi
    =
    \psi(\phi)-\mu\psi(A)-(1-\mu)\psi(B)
    +\mu\log y+(1-\mu)\log(1-y),
\]
\[
    L_{\phi\phi}
    =
    \psi_1(\phi)-\mu^2\psi_1(A)-(1-\mu)^2\psi_1(B).
\]
Then
\[
    a_\mu=\phi G,\qquad
    b_{\mu\mu}=-\phi^2\{\psi_1(A)+\psi_1(B)\},
\]
\[
    a_a=\phi L_\phi,\qquad
    B_{aa}=\phi L_\phi+\phi^2L_{\phi\phi},
\]
\[
    b_{\mu a}
    =
    \phi G+\phi^2\{-\mu\psi_1(A)+(1-\mu)\psi_1(B)\}.
\]

\subsubsection*{Zero-inflated Poisson}

Let $\zeta=\logit(\theta)$, $\theta=\{1+\exp(-\zeta)\}^{-1}$, and
$\dot\theta=\theta(1-\theta)$.  For $y>0$,
\[
    a_\mu=\frac{y}{\mu}-1,\qquad
    b_{\mu\mu}=-\frac{y}{\mu^2},\qquad
    a_\zeta=-\theta,\qquad
    B_{\zeta\zeta}=-\dot\theta,\qquad
    b_{\mu\zeta}=0 .
\]
For $y=0$, let $e=\exp(-\mu)$,
\[
    A=\theta+(1-\theta)e,\qquad
    B=(1-\theta)e,\qquad
    C=\dot\theta(1-e),
\]
and
\[
    C_\zeta=\dot\theta(1-2\theta)(1-e).
\]
Then
\[
    a_\mu=-\frac{B}{A},\qquad
    b_{\mu\mu}=\frac{B\theta}{A^2},
\]
\[
    a_\zeta=\frac{C}{A},\qquad
    B_{\zeta\zeta}=\frac{C_\zeta A-C^2}{A^2},
\]
\[
    b_{\mu\zeta}
    =
    \frac{\dot\theta e\,A+CB}{A^2}.
\]

\subsubsection*{Zero-inflated negative binomial}

For ZINB, use the NB2 block above for the non-zero count density and let
$\zeta=\logit(\theta)$.  For $y>0$, the derivatives in $(\mu,\rho)$ are exactly
the NB2 derivatives, while
\[
    a_\zeta=-\theta,\qquad B_{\zeta\zeta}=-\theta(1-\theta),
\]
with zero cross-derivatives involving $\zeta$.

For $y=0$, let $p_0$ be the NB2 probability of zero,
\[
    p_0=\left(\frac{r}{r+\mu}\right)^r,\qquad r=\exp(-\rho),
\]
and define
\[
    A=\theta+(1-\theta)p_0,\qquad
    C=\dot\theta(1-p_0),\qquad
    C_\zeta=\dot\theta(1-2\theta)(1-p_0).
\]
Let $g_j$ and $H_{j\ell}$ denote the NB2 score and Hessian entries for
$\log p_0$ in the coordinates $j,\ell\in\{\mu,\rho\}$.  Then
\[
    A_j=(1-\theta)p_0g_j,\qquad
    A_{j\ell}=(1-\theta)p_0(H_{j\ell}+g_jg_\ell),
\]
\[
    A_{j\zeta}=-\dot\theta p_0 g_j.
\]
The zero-count ZINB block is
\[
    a_j=\frac{A_j}{A},\qquad
    B_{j\ell}=\frac{A_{j\ell}A-A_jA_\ell}{A^2},
\]
\[
    a_\zeta=\frac{C}{A},\qquad
    B_{\zeta\zeta}=\frac{C_\zeta A-C^2}{A^2},
\qquad
    B_{j\zeta}=\frac{A_{j\zeta}A-A_jC}{A^2}.
\]

\subsubsection*{Skew normal}

For the Azzalini skew-normal component, let $\omega=\exp(a)$ be the scale,
$\alpha$ the shape, $z=(y-\mu)/\omega$, and $t=\alpha z$.  Let
$\varphi$ and $\Phi$ denote the standard normal density and distribution
function, and define the inverse Mills ratio
\[
    M(t)=\frac{\varphi(t)}{\Phi(t)},\qquad
    M'(t)=-M(t)\{t+M(t)\}.
\]
Write
\[
    \ell_z=-z+\alpha M(t),\qquad
    \ell_{zz}=-1+\alpha^2M'(t),\qquad
    \ell_{z\alpha}=M(t)+\alpha z M'(t).
\]
Then
\[
    a_\mu=-\frac{\ell_z}{\omega},\qquad
    b_{\mu\mu}=\frac{\ell_{zz}}{\omega^2},
\]
\[
    a_a=-1-z\ell_z,\qquad
    a_\alpha=zM(t),
\]
\[
    B_{aa}=z\ell_z+z^2\ell_{zz},\qquad
    B_{\alpha\alpha}=z^2M'(t),\qquad
    B_{a\alpha}=-z\ell_{z\alpha},
\]
\[
    b_{\mu a}=\frac{\ell_z+z\ell_{zz}}{\omega},\qquad
    b_{\mu\alpha}=-\frac{\ell_{z\alpha}}{\omega}.
\]

\begin{remark}[Proof of the family blocks]
Each displayed block is obtained by differentiating the corresponding
one-observation log density in $(\mu,\gamma)$ coordinates.  Positive nuisance
parameters are differentiated after the log transformation, so if
$u=\exp(\gamma)$ then
$\partial/\partial\gamma=u\,\partial/\partial u$ and
$\partial^2/\partial\gamma^2
=u\,\partial/\partial u+u^2\,\partial^2/\partial u^2$.  The ZIP and ZINB
zero-count formulas follow by differentiating the log of the zero mixture
probability, whereas the positive-count formulas add only the
$\log(1-\theta)$ contribution to the ordinary count-family likelihood.
\end{remark}

\subsection{Implementation validation}

The simulation code uses the analytic formulas above whenever available and
falls back to central finite differences only for families without analytic
blocks.  The validation script
\[
    \texttt{experiments/simulations/validate\_louis\_derivatives.py}
\]
compares each analytic block against central finite differences in
$(\eta,\gamma)$ coordinates.  This numerical check is not a mathematical proof,
but it is a useful guard against transcription and implementation errors after
the symbolic derivations have been fixed.

\section{Simulation Design Details}
\label{sec:supp_sim_design}

This section records the simulation settings used for the results in the main
manuscript.  The final simulation run was
\texttt{v4\_positive\_repaired\_louis}, with 100 Monte Carlo replications.
Scenario A and Scenario C used sample sizes
\[
    n\in\{500,1000,1500\},
\]
whereas Scenario B used \(n=1000\) and \(p=20\).  All training datasets were
paired with an independently generated test set of size 2000 for held-out
log-likelihood evaluation.

\subsection{Covariate generation and mixture sampling}

Each design matrix has an intercept in the first column.  The remaining
covariates are generated from a centered multivariate normal distribution with
AR(1) covariance
\[
    \Sigma_{ab}=0.2^{|a-b|}.
\]
Thus, for \(p=5\) there are four non-intercept covariates and for \(p=20\)
there are nineteen non-intercept covariates.  The data-generating covariates are
not centered or standardized before simulation.  The fitted models use internal
standardization for numerical stability; variable-selection summaries are
reported on the original coefficient scale, while the Wald coverage calculations
compare estimates and true coefficients on the fitted standardized scale.

For each observation, a latent label \(Z_i\) is sampled from
\(\Pr(Z_i=k)=\pi_k\).  Conditional on \(Z_i=k\), the linear predictor is
\(\eta_{ik}=x_i^\top\beta_k\), the component location or mean parameter is
\(\mu_{ik}=h_k(\eta_{ik})\), and \(Y_i\) is sampled from the corresponding
component family.  The four data-generating examples used in Scenarios A and C
are shown in Table~\ref{tab:supp_dgm_main}.

\begin{table}[htbp]
\centering
\caption{Data-generating models for Scenarios A and C.  The vectors displayed
are the raw regression coefficients, including the intercept.}
\label{tab:supp_dgm_main}
\resizebox{\textwidth}{!}{%
\begin{tabular}{lllll}
\toprule
Example & \(\pi\) & Component 1 & Component 2 & Notes \\
\midrule
1 & \((0.60,0.40)\) &
Gaussian, identity link, \(\beta_1=(1.0,0.5,-0.7,1.2,-0.4)\), \(\sigma=1\) &
Student-\(t\), identity link, \(\beta_2=(-0.5,-0.8,0.9,-1.0,0.6)\), df \(=4\), scale \(=1\) &
Real-valued heavy-tailed mixture \\
2 & \((0.55,0.45)\) &
Gamma, log link, \(\beta_1=(-0.60,0.20,-0.15,0.15,-0.10)\), shape \(=30\) &
Lognormal, identity link, \(\beta_2=(1.00,-0.35,0.30,-0.25,0.20)\), log-sd \(=1\) &
Positive-support mixture \\
3 & \((0.60,0.40)\) &
Poisson, log link, \(\beta_1=(2.5,0.3,-0.3,0.3,-0.3)\) &
NB2, log link, \(\beta_2=(0.0,-0.3,0.3,-0.3,0.3)\), \(\alpha=2\) &
Count mixture with overdispersion \\
4 & \((0.60,0.40)\) &
Poisson, log link, \(\beta_1=(2.0,0.3,-0.3,0.3,-0.3)\) &
ZIP, log link, \(\beta_2=(1.0,-0.3,0.3,-0.3,0.3)\), zero-inflation \(\theta=0.6\) &
Count mixture with excess zeros \\
\bottomrule
\end{tabular}}
\end{table}

For the Gamma component, the mean-shape parameterization is used:
\(\E(Y\mid Z=k,X)=\mu\), \(\operatorname{Var}(Y\mid Z=k,X)=\mu^2/\text{shape}\).
For NB2, the variance is \(\mu+\alpha\mu^2\).  For the lognormal component,
\(\mu\) denotes the log-location; the response mean is
\(\exp(\mu+\sigma^2/2)\).  For ZIP, the response mean is
\((1-\theta)\mu\).

\subsection{Scenario A: model selection and prediction}

Scenario A evaluates family-tuple recovery and held-out prediction.  For each
example and sample size, the fitted candidate space is restricted to families
with the same support as the data-generating example:
\[
\begin{array}{ll}
\text{real support:} & \{\text{Gaussian},\text{Student-}t,\text{skew normal}\},\\
\text{positive support:} & \{\text{Gamma},\text{lognormal},\text{inverse Gaussian}\},\\
\text{count support:} & \{\text{Poisson},\text{NB2},\text{ZIP}\}.
\end{array}
\]
The search uses \(K_{\max}=3\), Bayesian information criterion (BIC) ranking,
beam width 10, two random starts per candidate, maximum 100
expectation-maximization (EM) iterations, tolerance \(10^{-4}\), and internal
standardization.  Initializations are quantile-based for real-valued responses,
\(Y\)-clustered GLM initialization for positive responses, and random
initialization for count responses.

The oracle model is the fitted model with the true two-component family tuple.
The selected model is the lowest-BIC candidate from the beam search.  Prediction
is measured by average held-out log likelihood on the independent test set.
For each replication, the ten best models by BIC and the ten best models by
held-out log likelihood are retained to form the leaderboard summaries.

Strict recovery requires the selected family tuple to match the true tuple up to
label permutation.  Boundary-equivalent recovery is reported separately using
the rules in Table~\ref{tab:supp_equiv}.  These rules are used only for
interpretation of families with nesting or boundary behavior; strict recovery is
always reported separately in the main table.

\begin{table}[htbp]
\centering
\caption{Equivalence classes used in Scenario A summaries.}
\label{tab:supp_equiv}
\begin{tabular}{ll}
\toprule
True unordered tuple & Treated as equivalent for summary purposes \\
\midrule
\(\{\text{Gaussian},\text{Student-}t\}\) &
\(\{\text{Gaussian},\text{Student-}t\}\), \(\{\text{Student-}t,\text{Student-}t\}\) \\
\(\{\text{Gamma},\text{lognormal}\}\) &
\(\{\text{Gamma},\text{lognormal}\}\) \\
\(\{\text{Poisson},\text{NB2}\}\) &
\(\{\text{Poisson},\text{NB2}\}\), \(\{\text{NB2},\text{NB2}\}\) \\
\(\{\text{Poisson},\text{ZIP}\}\) &
\(\{\text{Poisson},\text{ZIP}\}\), \(\{\text{ZIP},\text{ZIP}\}\) \\
\bottomrule
\end{tabular}
\end{table}

\subsection{Scenario B: component-specific variable selection}

Scenario B evaluates sparse component-specific regression recovery with the
true family tuple fixed.  It uses Examples 1, 2, and 4, with \(n=1000\) and
\(p=20\).  The first five entries of each coefficient vector are nonzero
including the intercept, and the remaining fifteen entries are zero.  Let
\(\mathbf 0_{15}\) denote fifteen trailing zeros.  The coefficients are shown
in Table~\ref{tab:supp_sparse_betas}.

\begin{table}[htbp]
\centering
\caption{Sparse coefficient vectors used in Scenario B.}
\label{tab:supp_sparse_betas}
\resizebox{\textwidth}{!}{%
\begin{tabular}{lll}
\toprule
Example & Component 1 coefficient vector & Component 2 coefficient vector \\
\midrule
1 &
\((1.0,0.8,-0.6,1.2,-0.5,\mathbf 0_{15})\) &
\((-0.8,-1.0,0.9,-1.2,0.7,\mathbf 0_{15})\) \\
2 &
\((-0.60,0.20,-0.15,0.15,-0.10,\mathbf 0_{15})\) &
\((1.00,-0.35,0.30,-0.25,0.20,\mathbf 0_{15})\) \\
4 &
\((2.0,0.4,-0.3,0.6,-0.25,\mathbf 0_{15})\) &
\((1.0,-0.5,0.45,-0.6,0.35,\mathbf 0_{15})\) \\
\bottomrule
\end{tabular}}
\end{table}

For each replication, three penalty classes are compared: no penalty, lasso,
and elastic net.  The lasso and elastic-net tuning grid is
\[
    \lambda\in\{0.25,0.5,1,2,5,10,20,50\},
\]
and the elastic-net mixing parameter is fixed at \(0.5\).  Within each penalty
class, the representative model reported in the main table is the member of the
path with the lowest BIC.  A coefficient is counted as selected when its
absolute value exceeds \(10^{-4}\).  The true-positive rate (TPR) and
false-positive rate (FPR) are averaged over the two components, and the
false-discovery rate (FDR), false-negative rate (FNR), coefficient mean squared
error (MSE), selected model size, and held-out log likelihood are recorded for
the selected member of each penalty class.

\subsection{Scenario C: standard errors and coverage}

Scenario C evaluates uncertainty quantification for the unpenalized estimator
with the true family tuple fixed.  It uses Examples 1, 2, and 4 at
\(n\in\{500,1000,1500\}\).  For each fit, two observed-information calculations
are compared:
\[
    \text{analytic Louis information}
    \qquad\text{and}\qquad
    \text{central finite-difference Hessian}.
\]
The Louis calculation uses the analytic derivative blocks in Section~S1 when
available and finite differences only for families without analytic blocks.
The reported intervals are nominal 95\% Wald intervals.  Because the fitting
routine standardizes the design internally, coverage for regression
coefficients is computed after mapping the true raw coefficients into the
standardized fitting scale and after aligning component labels by the nearest
coefficient vector.

\subsection{Reproducibility settings}

The simulation seeds are deterministic.  For replication index
\(r=0,\ldots,99\) and example \(e\), the seed is \(100r+e\).  The same seeded
random-number stream is used to generate the training sample and then the
independent test sample.  All scenarios use a maximum of 100 EM iterations,
tolerance \(10^{-4}\), two starts, and internal standardization in the fitted
models.  The repaired simulation output used in the main manuscript is stored
under
\begin{quote}
\footnotesize
\texttt{experiments/simulations/paper\_outputs/}\\
\texttt{v4\_positive\_repaired\_20260624\_b}.
\end{quote}

\section{Application Details and Additional Results}
\label{sec:supp_applications}

\subsection{Parkinson's telemonitoring}

The Parkinson's telemonitoring data were obtained from the UCI Machine Learning
Repository \citep{uci2009parkinsons,tsanas2010parkinsons}.  The source contains
5,875 recordings from 42 participants.  We used
\(Y=\log(\mathtt{total\_UPDRS})\), where UPDRS denotes the Unified Parkinson's
Disease Rating Scale.  We excluded the participant identifier and motor UPDRS
from the design, and retained age, sex, test time, and the 16
available voice measurements.  There were no nonfinite values after
preparation.

The computational screen sampled 3,000 recordings without replacement, used
2,000 for training and 1,000 for testing, and screened at most 40 predictors
using the training sample only.  The candidate family set was Gaussian,
Student-$t$, and skew normal.  All unordered tuples with
\(K\in\{1,2,3\}\) were crossed with
\[
  \lambda\in\{0,0.1,0.25,0.5,1,2,5,10,20\}
\]
and two initialization schemes.  Fits were allowed 160 EM iterations with
tolerance \(10^{-3}\) and two starts.  Of 342 candidates, 319 satisfied the
recorded convergence criterion.  Only converged fits were eligible for the
reported leaderboards.

Table~\ref{tab:supp_park_bic} gives the leading converged family structures,
after retaining the best tuning and initialization for each tuple.  A
Student-$t$--Student-$t$--Student-$t$ fit at \(\lambda=20\) had provisional
BIC 1417.9 but reached the 160-iteration limit and was excluded.  The converged
fit for that family tuple at \(\lambda=10\) had BIC 1496.4.  This diagnostic
is reported because mixture rankings should not silently treat an unfinished
optimization as a valid likelihood maximum.

\begin{table}[htbp]
\centering
\caption{Leading converged Parkinson's family structures by BIC.  SN denotes
skew normal, and RMSE denotes root mean squared error.  Scores are based on
2,000 training and 1,000 held-out recordings.}
\label{tab:supp_park_bic}
\small
\begin{tabular}{lrrrr}
\toprule
Family tuple & \(\lambda\) & BIC & Test log score & RMSE \\
\midrule
Student-\(t\)+SN+SN & 20 & 1452.2 & -297.3 & .392 \\
Student-\(t\)+Student-\(t\)+SN & 20 & 1475.3 & -302.6 & .393 \\
Gaussian+Gaussian+Gaussian & 20 & 1495.6 & -287.3 & .385 \\
Student-\(t\)+Student-\(t\)+Student-\(t\) & 10 & 1496.4 & -334.5 & .391 \\
Gaussian+Gaussian+SN & 20 & 1503.0 & -286.8 & .385 \\
Gaussian+Student-\(t\)+SN & 20 & 1510.0 & -288.3 & .385 \\
Gaussian+Gaussian+Student-\(t\) & 20 & 1526.1 & -287.1 & .385 \\
Gaussian+SN+SN & 20 & 1526.2 & -287.4 & .385 \\
SN+SN+SN & 20 & 1533.6 & -287.2 & .385 \\
Gaussian+Student-\(t\)+Student-\(t\) & 20 & 1552.1 & -350.0 & .385 \\
\bottomrule
\end{tabular}
\end{table}

The BIC-selected model had active-set sizes 14, 12, and 16.  Age, sex, test
time, Jitter(\%), Jitter(Abs), Shimmer:APQ11, NHR, HNR, RPDE, DFA, and PPE
were shared.  Component 1 additionally retained Jitter:PPQ5,
Shimmer:APQ3, and Shimmer:APQ5.  Component 2 additionally retained Shimmer.
Component 3 additionally retained Jitter:RAP, Jitter:PPQ5, Jitter:DDP,
Shimmer, and Shimmer:APQ3.  These supports are descriptive results from one
recording-level split; no post-selection coefficient-level inference is
claimed.

Because the same participant can contribute recordings to both partitions, the
test scores do not estimate performance for unseen participants.  A
participant-grouped analysis is required before making that stronger
generalization.

\subsection{RAND baseline cross-section}

The RAND data were obtained from the \texttt{randhealth} data in the
\texttt{camerondata} collection, which derives from the RAND Health Insurance
Experiment \citep{deb2002health,cameron2005microeconometrics}.  The source has
20,190 participant-year records.  We sorted by participant and study year and
retained the earliest record, leaving 5,912 independent participant records.
The response is the number of non-physician visits.  Outcome variables and the
participant identifier were excluded from the predictors; plan, site, and year
were represented by indicators.

For split validation, three independent 80/20 participant partitions were
used.  Screening was repeated within each training partition and retained at
most 50 predictors.  The search crossed all unordered tuples from Poisson,
NB2, ZIP, and ZINB for \(K=1,2,3\), ten penalty values, and two
initializations, yielding 1,360 fits per split and 4,080 fits in total.  The
Poisson--NB2 tuple won BIC in all three splits.  Its held-out average log scores
were \(-0.642\), \(-0.628\), and \(-0.683\).  The corresponding best
identical-mixture values were \(-0.641\), \(-0.628\), and \(-0.682\), so
the family choice is not a claim of superior point or log-score prediction over
an NB2--NB2 approximation.

The full-data search used the same family and tuning spaces and produced 680
fits.  All were recorded as converged.  Table~\ref{tab:supp_rand_full} lists
the leading distinct family structures.  The best NB2--NB2 model estimated
\(\log\alpha_1=-8.518\), while the NB2 component of the selected
Poisson--NB2 model had \(\log\alpha_2=2.345\).

\begin{table}[htbp]
\centering
\caption{Leading distinct family structures in the full-data RAND search.}
\label{tab:supp_rand_full}
\small
\begin{tabular}{lrrrr}
\toprule
Family tuple & \(K\) & \(\lambda\) & BIC & Active counts \\
\midrule
Poisson+NB2 & 2 & 100 & 7950.6 & 6/0 \\
Poisson+ZINB & 2 & 50 & 7954.4 & 16/0 \\
NB2+NB2 & 2 & 100 & 7955.2 & 6/0 \\
NB2+ZINB & 2 & 100 & 7963.8 & 6/0 \\
Poisson+Poisson+NB2 & 3 & 100 & 7970.1 & 6/0/0 \\
ZINB+ZINB & 2 & 100 & 7970.4 & 7/0 \\
Poisson+Poisson+ZINB & 3 & 100 & 7970.5 & 6/0/0 \\
Poisson+NB2+ZINB & 3 & 100 & 7972.3 & 6/0/0 \\
NB2 & 1 & 50 & 8504.9 & 18 \\
\bottomrule
\end{tabular}
\end{table}

For the final family pair, an unpenalized constrained refit was performed on
the support selected at \(\lambda=100\).  The analytical Louis matrix was
reduced to the ten free coordinates: one mixing logit, two intercepts, six
Poisson slopes, and NB2 log dispersion.  It had rank 10, condition number
1439.2, and minimum eigenvalue 1.280.  Every observation used analytical
Poisson and NB2 derivative blocks.

The participant bootstrap held the family pair fixed but repeated screening,
tuning, support selection, and refitting.  The selected penalties and
component support sizes are shown in Table~\ref{tab:supp_rand_bootdiag}.  The
NB2 component was intercept-only in 162 of 500 samples.  Its most frequently
selected slopes were physical limitation (0.450), child status (0.384), site 6
(0.360), disease burden (0.346), and age (0.316); none exceeded 0.5.  These
results motivate the cautious support interpretation in the main text.

\begin{table}[htbp]
\centering
\caption{RAND bootstrap tuning and support diagnostics.}
\label{tab:supp_rand_bootdiag}
\small
\begin{tabular}{lrrr}
\toprule
Quantity & Lower 2.5\% & Median & Upper 97.5\% \\
\midrule
Poisson active-set size & 5 & 16 & 27.5 \\
NB2 active-set size & 0 & 2 & 13 \\
Poisson weight & .809 & .835 & .868 \\
NB2 dispersion & 4.330 & 8.049 & 12.159 \\
\midrule
Selected \(\lambda\) & \multicolumn{3}{l}{25: 13.0\%; 50: 67.4\%; 100: 19.6\%} \\
\bottomrule
\end{tabular}
\end{table}

\bibliographystyle{plainnat}
\bibliography{references}

\end{document}
